\section{Introduction}\label{sec:intro}
An experiment does not supply a geometric state space to its observer. It supplies preparations, admissible interventions and reports. A state is predictive only to the extent that it determines executable future tests; a geometric resolution is useful only to the extent that it is acquired with positive probability and can be retained by a causal machine. The problem is to pass from these primitive data to a sharp resource law without treating a formal reachable set, a free exact posterior, or a phase-dependent program as an acquired finite resource.

There are two distinct obstacles. Static quantization may discard information needed to choose subsequent observations. Even when a compressed state predicts the next report well, the continuation value of choosing the wrong action can be nonsmooth. A pointwise Lipschitz estimate then loses an order relative to the quadratic terminal task. Moreover, stability of a passive filter does not by itself compare a finite-state feedback controller with the optimal full-history controller of the same experiment.

The main theorem resolves these obstacles on a specified, nontrivial class of controlled positive instruments. Its new stability mechanism is an overlap of the \emph{actual acquired laws}, uniform over feedback choices. If every controlled block retains a common positive component, the oscillation of every terminal Bellman continuation decays. Concavity and an interior preparation buffer convert this oscillation bound into a local Lipschitz bound. An acquired-density estimate then integrates the nonsmooth supporting-plane error to second order. The resulting controller chooses a Bellman-optimal action at its own calibrated retained state, not at an inaccessible full-history posterior.

Write $V_0(u_0)$ for the unrestricted full-history adaptive optimum of the declared terminal task. Under the hypotheses of Theorem~\ref{thm:main}, for every $n\ge1$,
\begin{equation}\label{eq:intro-law}
 \calR_{\rm cp}(n,m)-V_0(u_0)\asymp m^{-2/d},\qquad
 \calR_{\rm aut}(n,M)-V_0(u_0)\asymp(M-n)^{-2/d}.
\end{equation}
The second class is nonempty precisely when $M\ge n+1$. Its $n$-state charge is an internal exact-deadline charge, not a change from labels to bits. The constants are uniform in both budgets for fixed experiment certificates. They need not remain bounded as acquired overlap disappears. The baseline concerns the same dynamic terminal experiment for each $n$; no monotone static-parameter learning curve is inferred.

Two raw-kernel consequences give the theorem operational content. In the classical experiment, an observer selects among Gaussian sensors before each report, while the hidden state subsequently evolves through any fixed positive invertible transition. The theorem includes sensor steps whose acquired belief kernel expands Euclidean Wasserstein distance. In the quantum experiment, a report-dependent posterior is followed by a chosen unitary and a strictly positive channel. Small-effect report corridors yield feedback-uniform overlap even when one-step transport estimates do not contract. An asymmetric informationally complete detector makes the feedback choice strictly depend on the current acquired state. These results concern known models, finite action sets, and a finite terminal task; they do not assert unknown-kernel learning, arbitrary stopping, or universal control optimality beyond this interface.

The conditional-calibration identity, its Bellman--Jensen form, the integrated convex-curvature estimate, and the internal-clock argument have antecedents in the retained work \cite{GTFControl,GTFAllocation}. The likelihood-normalized three-defect identity was established separately in \cite{GTFBayes}. We reprove the needed results here so that the native article is self-contained. The new connection is the route from feedback-uniform \emph{acquired overlap} to the sharp controlled total-state law, together with its active Gaussian, feedback quantum and rank-changing verifications. Neither a conditional projection identity nor the classical common-mass contraction principle is claimed as a new general fact.

The distinction from value-directed belief compression \cite{PB2000,PB2001,PB2002}, finite-memory feedback approximation \cite{KY2022}, and recent belief-window and partition-filter approximations \cite{Kim2026,DGM2026} is the joint guarantee: an implemented autonomous controller is compared with the same full-history optimum, while a converse applies to every randomized Borel controller with the stated total alphabet. Section~\ref{sec:comparison} makes the comparison at the level of tasks, memory objects and theorem hypotheses. Section~\ref{sec:resources} treats numerical execution and causal simulation separately from hard-state existence.
